\documentclass[ ../main.tex]{subfiles}
\providecommand{\mainx}{..}
\begin{document}
\section{The Bernoulli map model}
\label{sec:bernoulli-map}

The set of total functions of type $\Set{X} \to \Set{Y}$ is denoted by the
exponential $\Set{Y}^{\Set{X}}$ and has cardinality
$\Card{\Set{Y}}^{\Card{\Set{X}}}$.
Given a fixed function $\Fun{f} \colon \Set{X} \to \Set{Y}$, a function
$\Fun{g} \colon \Set{X} \to \Set{Y}$ chosen uniformly at random from
$\Set{Y}^{\Set{X}}$ agrees with $\Fun{f}$ at a given point $x \in \Set{X}$
with probability $\Card{\Set{Y}}^{-1}$, and agrees everywhere---that is,
$\Fun{g} = \Fun{f}$---with probability $\Card{\Set{Y}}^{-\Card{\Set{X}}}$.

We seek a more refined model.
Rather than selecting an approximation uniformly at random from the full
function space, we model approximation errors that arise from rate distortion,
noise, or incomplete knowledge as \emph{per-element Bernoulli trials}~\cite{feller1968} with
structured error rates.
This is the \emph{Bernoulli map model}, which we develop in this section.

%% ------------------------------------------------------------------
\subsection{The general model}
\label{sec:general_model}

Let $\Fun{f} \colon \Set{X} \to \Set{Y}$ be an objective (exact) function and
let $\APFun{f} \colon \Set{X} \to \Set{Y}$ be a random function that
approximates $\Fun{f}$.
We measure pointwise correctness with the \emph{Iverson bracket}.

\begin{definition}[Iverson bracket]
\label{def:iverson}
The Iverson bracket $\llbracket \cdot \rrbracket \colon \{\True, \False\} \to
\BitSet$ maps a predicate to $1$ if the predicate is true and to $0$ otherwise.
\end{definition}

For each $x \in \Set{X}$, define the \emph{correctness indicator}
\begin{equation}
\label{eq:correctness}
    \RV{Y}_x \;\coloneqq\; \llbracket \Fun{f}(x) = \APFun{f}(x) \rrbracket\,.
\end{equation}
Since $\APFun{f}$ is random, $\RV{Y}_x$ is a Bernoulli random variable taking
values in $\BitSet$.
The probability that the approximation is correct on the entire domain is
\begin{equation}
\label{eq:joint_correct}
    \Prob{\bigcap_{x \in \Set{X}} \{\RV{Y}_x = 1\}}
    \;=\; \prod_{x \in \Set{X}} \Prob{\RV{Y}_x = 1}\,,
\end{equation}
provided the family $\{\RV{Y}_x\}_{x \in \Set{X}}$ consists of mutually
independent random variables.
This independence assumption is fundamental to the model and we state it
explicitly.

\begin{assumption}[Element-wise independence]
\label{asm:independence}
The correctness indicators $\{\RV{Y}_x\}_{x \in \Set{X}}$ are mutually
independent.
\end{assumption}

Under \cref{asm:independence}, each $\RV{Y}_x$ is an independent Bernoulli
random variable with some parameter $p_x \coloneqq \Prob{\RV{Y}_x = 1}$.
In general, $p_x$ may differ across elements of $\Set{X}$.
The \emph{order} of the model captures the complexity of this heterogeneity.

\begin{definition}[Model order]
\label{def:model_order}
Let $\APFun{f}$ be a random approximate map of $\Fun{f} \colon \Set{X} \to
\Set{Y}$ satisfying \cref{asm:independence}.
The \emph{order} of the random approximate map is the minimum integer $k \geq 0$
such that there exists a partition of $\Set{X}$ into $k$ non-empty blocks
$\Set{X}_1, \ldots, \Set{X}_k$ with the property that within each block
$\Set{X}_i$, the correctness indicators $\{\RV{Y}_x : x \in \Set{X}_i\}$ are
independent and identically distributed as $\berdist(p_i)$ for some
$p_i \in [0,1]$.
\end{definition}

\begin{remark}
A \emph{zeroth-order} model ($k=0$) corresponds to the degenerate case of an
exact function: $\RV{Y}_x = 1$ with probability $1$ for all $x \in \Set{X}$,
so $\APFun{f} = \Fun{f}$ almost surely.
A \emph{first-order} model ($k=1$) posits a single partition block comprising
all of $\Set{X}$, so that the indicators $\{\RV{Y}_x\}_{x \in \Set{X}}$ are
i.i.d.\ $\berdist(p)$.
A \emph{second-order} model ($k=2$) partitions the domain into two blocks with
distinct Bernoulli parameters; this is the natural setting for false positive
and false negative rates, which we develop in the next subsection.
\end{remark}

Since within each partition block $\Set{X}_i$ the indicators are i.i.d.\
$\berdist(p_i)$, the number of correct outputs in block $\Set{X}_i$ is
binomially distributed:
\begin{equation}
\label{eq:binom_block}
    \sum_{x \in \Set{X}_i} \RV{Y}_x
    \;\sim\; \bindist\!\left(\Card{\Set{X}_i},\, p_i\right)\,.
\end{equation}

%% ------------------------------------------------------------------
\subsection{The second-order model}
\label{sec:second_order}

The most common instantiation of the Bernoulli map model is the second-order
model, in which the domain $\Set{X}$ is partitioned into two blocks determined
by the structure of the objective function $\Fun{f}$.

Let $\Fun{f} \colon \Set{X} \to \Maybe{\Set{Y}}$ be a partial function with
domain of definition $\PlainSet{P} \coloneqq \Dod(\Fun{f}) \subseteq \Set{X}$
(the \emph{positives}) and let $\PlainSet{N} \coloneqq \Set{X} \smallsetminus
\PlainSet{P}$ (the \emph{negatives}).
The set-indicator function of the domain of definition partitions $\Set{X}$
into these two blocks:
\begin{equation}
\label{eq:partition}
    \Set{X} = \PlainSet{P} \;\dot{\cup}\; \PlainSet{N}\,,
    \qquad
    \PlainSet{P} \cap \PlainSet{N} = \EmptySet\,.
\end{equation}

Within each block, the correctness indicators are i.i.d., but with
\emph{different} Bernoulli parameters.
The error rates on these two blocks are the false positive rate and the false
negative rate, defined as conditional losses with respect to the framework
introduced in \cref{sec:error}.

\begin{definition}[False positive rate]
\label{def:fpr}
Let $\Fun{f} \colon \Set{X} \to \Maybe{\Set{Y}}$ be a partial function with
negatives $\PlainSet{N} = \Set{X} \smallsetminus \Dod(\Fun{f})$, and let
$\APFun{f}$ be a random approximate map of $\Fun{f}$.
The \emph{false positive rate} is the conditional expected loss over
$\PlainSet{N}$:
\begin{equation}
\label{eq:fpr_def}
    \fprate
    \;\coloneqq\;
    \Fun{\ell}_{\RV{X} \Given \PlainSet{N}}\!\left(\APFun{f} \Given \Fun{f}\right)
    \;=\;
    \Prob{\APFun{f}(\RV{X}) \neq \Nothing \Given \RV{X} \in \PlainSet{N}}\,.
\end{equation}
The complementary quantity $\tnrate \coloneqq 1 - \fprate$ is the \emph{true
negative rate}.
\end{definition}

\begin{definition}[False negative rate]
\label{def:fnr}
Let $\Fun{f} \colon \Set{X} \to \Maybe{\Set{Y}}$ be a partial function with
positives $\PlainSet{P} = \Dod(\Fun{f})$, and let $\APFun{f}$ be a random
approximate map of $\Fun{f}$.
The \emph{false negative rate} is the conditional expected loss over
$\PlainSet{P}$:
\begin{equation}
\label{eq:fnr_def}
    \fnrate
    \;\coloneqq\;
    \Fun{\ell}_{\RV{X} \Given \PlainSet{P}}\!\left(\APFun{f} \Given \Fun{f}\right)
    \;=\;
    \Prob{\APFun{f}(\RV{X}) \neq \Fun{f}(\RV{X}) \Given \RV{X} \in \PlainSet{P}}\,.
\end{equation}
The complementary quantity $\tprate \coloneqq 1 - \fnrate$ is the \emph{true
positive rate}.
\end{definition}

The total expected loss over $\Set{X}$ decomposes as a weighted sum of these
conditional losses.
By the partition $\Set{X} = \PlainSet{P} \,\dot{\cup}\, \PlainSet{N}$ and
the law of total expectation (cf.\ \cref{eq:disj_loss}), we obtain:

\begin{lemma}[Loss decomposition]
\label{prop:loss_decomp}
Let $\RV{X}$ be uniformly distributed on $\Set{X}$.
Then the total expected loss of $\APFun{f}$ with respect to $\Fun{f}$ is
\begin{equation}
\label{eq:loss_decomp}
    \Fun{\ell}_{\RV{X}}\!\left(\APFun{f} \Given \Fun{f}\right)
    \;=\;
    \fprate \cdot \frac{\Card{\PlainSet{N}}}{\Card{\Set{X}}}
    \;+\;
    \fnrate \cdot \frac{\Card{\PlainSet{P}}}{\Card{\Set{X}}}\,.
\end{equation}
\end{lemma}
\begin{proof}
By the law of total probability applied to the partition
$\Set{X} = \PlainSet{P} \,\dot{\cup}\, \PlainSet{N}$,
\begin{align}
    \Fun{\ell}_{\RV{X}}\!\left(\APFun{f} \Given \Fun{f}\right)
    &= \Fun{\ell}_{\RV{X} \Given \PlainSet{N}}\!\left(\APFun{f} \Given \Fun{f}\right)
       \cdot \Prob{\RV{X} \in \PlainSet{N}}
     + \Fun{\ell}_{\RV{X} \Given \PlainSet{P}}\!\left(\APFun{f} \Given \Fun{f}\right)
       \cdot \Prob{\RV{X} \in \PlainSet{P}} \notag \\
    &= \fprate \cdot \frac{\Card{\PlainSet{N}}}{\Card{\Set{X}}}
     + \fnrate \cdot \frac{\Card{\PlainSet{P}}}{\Card{\Set{X}}}\,,
\end{align}
where the second equality uses \cref{def:fpr,def:fnr} and the uniformity of
$\RV{X}$.
\end{proof}

The Bernoulli trial structure of the second-order model yields binomial error
counts.

\begin{lemma}[Binomial error counts]
\label{prop:binom_errors}
Let $\APFun{f}[\fprate][\fnrate]$ be a second-order random approximate map of
$\Fun{f}$.
Then:
\begin{enumerate}
    \item Each element $x \in \PlainSet{N}$ independently generates a false
          positive with probability $\fprate$, and the total number of false
          positives is
          $\FP \;\sim\; \bindist\!\left(\Card{\PlainSet{N}},\, \fprate\right)$.
    \item Each element $x \in \PlainSet{P}$ independently generates a false
          negative with probability $\fnrate$, and the total number of false
          negatives is
          $\FN \;\sim\; \bindist\!\left(\Card{\PlainSet{P}},\, \fnrate\right)$.
    \item The random variables $\FP$ and $\FN$ are independent.
\end{enumerate}
\end{lemma}
\begin{proof}
By \cref{def:model_order}, within the block $\PlainSet{N}$ the correctness
indicators are i.i.d.\ $\berdist(1-\fprate)$.
The indicator for a false positive at $x \in \PlainSet{N}$ is
$1 - \RV{Y}_x \sim \berdist(\fprate)$, and the count
$\FP = \sum_{x \in \PlainSet{N}} (1 - \RV{Y}_x)$ is a sum of
$\Card{\PlainSet{N}}$ i.i.d.\ $\berdist(\fprate)$ trials, hence
$\FP \sim \bindist(\Card{\PlainSet{N}}, \fprate)$.
The argument for $\FN$ on $\PlainSet{P}$ is identical with parameter $\fnrate$.
Independence of $\FP$ and $\FN$ follows from the element-wise independence
(\cref{asm:independence}) and the disjointness of $\PlainSet{P}$ and
$\PlainSet{N}$.
\end{proof}

\begin{remark}
If the set $\PlainSet{P}$ is known a priori to both the producer and the
consumer of the approximate map, then queries may be restricted to
$\PlainSet{P}$ and the false positive rate $\fprate$ becomes irrelevant.
In many practical settings, however, $\PlainSet{P}$ is not known to the
querier, and $\fprate$ governs the noise on out-of-domain queries.
\end{remark}

%% ------------------------------------------------------------------
\subsection{First-order rate-distorted maps}
\label{sec:first_order}

We now describe a concrete generative model for the approximation errors.
When the approximation arises from \emph{rate distortion}---that is, from
encoding $\Fun{f}$ in a lossy compressed representation---the errors on
negatives have a specific distributional structure determined by the encoding
scheme.

\begin{definition}[First-order random approximate map]
\label{def:first_order_map}
Let $\Fun{f} \colon \Set{X} \to \Set{Y}$ be a function and let
$\Set{A} \subseteq \Set{X}$.
A \emph{first-order random approximate map} of $\Fun{f}$ with respect to
$\Set{A}$ is a random function
$\APFun{f} \colon \Set{X} \to \Set{Y}$ defined by
\begin{equation}
\label{eq:first_order_def}
    \APFun{f}(x) \;\coloneqq\;
    \begin{cases}
        \Fun{f}(x)     & \text{if } x \in \Set{A}\,,\\
        \RV{Y}^{+}     & \text{if } x \notin \Set{A}\,,
    \end{cases}
\end{equation}
where $\RV{Y}^{+}$ is a random variable with support contained in $\Set{Y}$,
drawn independently for each $x \notin \Set{A}$.
\end{definition}

In a first-order model, the output is \emph{deterministically correct} on
$\Set{A}$ (the false negative rate is zero) and \emph{random} on
$\Set{X} \smallsetminus \Set{A}$.
The distribution of the false positive value $\RV{Y}^{+}$ depends on the
encoding scheme used to represent $\Fun{f}$.

\begin{definition}[Rate-distorted false positive distribution]
\label{def:fp_distribution}
Let $\Encode_{\Set{Y}} \colon \Set{Y} \to \BitSet^{*}$ be a prefix-free
encoder for values in $\Set{Y}$.
The \emph{false positive value distribution} induced by $\Encode_{\Set{Y}}$ is
the probability mass function
\begin{equation}
\label{eq:fp_pmf}
    \Fun{p}_{\RV{Y}^{+}}(y)
    \;=\;
    \sum_{b \,\in\, \Encode_{\Set{Y}}(y)} 2^{-\BL(b)}\,,
\end{equation}
where $\BL(b)$ denotes the bit length of the code word $b$.
\end{definition}

The false positive value distribution arises because, for elements outside
$\Set{A}$, the data structure stores no information about $\Fun{f}(x)$.
A query on such an element produces a random bit string whose decoding under
$\Encode_{\Set{Y}}$ determines the spurious output value.
Shorter code words are decoded more frequently, so the distribution
\eqref{eq:fp_pmf} is naturally biased toward values with shorter
representations.

\begin{remark}
An optimal prefix-free encoder may be constructed by computing the frequency
of each element $y \in \Set{Y}$ in the image $\Fun{f}(\Set{A})$ and generating
a Huffman code.
However, storing a Huffman code table may require substantial memory when
$\Card{\Set{Y}}$ is large.
Practical alternatives include:
\begin{enumerate}
    \item prefix codes for known distributional families (e.g., Golomb codes
          for geometric distributions);
    \item universal codes (e.g., Elias codes) when $\Card{\Set{Y}}$ is very
          large or the distribution is unknown;
    \item uniform codes of length $\lceil \log_2 \Card{\Set{Y}} \rceil$, which
          yield $\Fun{p}_{\RV{Y}^{+}}(y) = \Card{\Set{Y}}^{-1}$ for all
          $y \in \Set{Y}$.
\end{enumerate}
The uniform encoder maximizes the entropy of $\RV{Y}^{+}$ and is the simplest
to implement, but is optimal only when $\Fun{f}$ is itself uniformly
distributed over $\Set{Y}$.
\end{remark}

The following example shows that the Bernoulli set---the central object of
study in the companion paper \cite{aset}---is a special case of the first-order
random approximate map.

\begin{example}
\label{ex:aset_as_map}
Let $\Set{A} \subseteq \Set{X}$ and consider the set-indicator function
$\SetIndicator{\Set{A}} \colon \Set{X} \to \BitSet$ defined by
$\SetIndicator{\Set{A}}(x) = \llbracket x \in \Set{A} \rrbracket$.
The first-order positive random approximate set $\PASet{A}[\fprate]$ is the
first-order random approximate map of $\SetIndicator{\Set{A}}$ with respect to
$\Set{A}$, where the false positive value distribution on $\BitSet = \{0,1\}$
satisfies
\begin{equation}
    \Fun{p}_{\RV{Y}^{+}}(1) = \fprate\,,
    \qquad
    \Fun{p}_{\RV{Y}^{+}}(0) = 1 - \fprate\,.
\end{equation}
Indeed, a false positive occurs when $x \notin \Set{A}$ yet the approximate map
outputs $1$ (indicating membership), which happens with probability $\fprate$.
Since $\Card{\BitSet} = 2$, a single bit suffices to encode the codomain, and
the false positive rate is entirely determined by the bias of that bit.
\end{example}

%% ------------------------------------------------------------------
\subsection{Notation}
\label{sec:notation}

We adopt the following notational conventions for random approximate maps
throughout this paper.

Let $\Fun{f} \colon \Set{X} \to \Maybe{\Set{Y}}$ be a partial function.
\begin{enumerate}
    \item The \emph{random approximate map} of $\Fun{f}$ with false positive
          rate $\fprate$ and false negative rate $\fnrate$ is denoted
          \begin{equation}
          \label{eq:apfun_notation}
              \APFun{f}[\fprate][\fnrate] \colon \Set{X} \to \Maybe{\Set{Y}}\,.
          \end{equation}
    \item When we wish to emphasize the \emph{type} of the approximate map
          rather than the function name, we write
          \begin{equation}
          \label{eq:arrow_notation}
              \Set{X} \amapsto{\fprate}{\fnrate} \Maybe{\Set{Y}}\,.
          \end{equation}
    \item When the error rates are unknown, irrelevant, or clear from context,
          we write simply $\APFun{f}$.
\end{enumerate}

\begin{remark}
The superscript/subscript convention on $\APFun{f}[\fprate][\fnrate]$ places
the false positive rate $\fprate$ in the superscript position and the false negative rate $\fnrate$ in the subscript position,
mirroring the notation $\ASet{A}[\fprate][\fnrate]$ for random approximate sets.
The arrow notation $\Set{X} \amapsto{\fprate}{\fnrate} \Set{Y}$ similarly
annotates the mapping arrow with the error rates, making the rates part of the
\emph{type signature} of the approximate map.
\end{remark}

\end{document}
